\subsection{Structural Relationship Evaluation}

The previous operations produce or mutate sets. Structural relationship methods answer a different kind of question. They do not ask for a new membership domain. They ask whether one set has a specific relationship to another set.

The main relationship questions are:

\begin{itemize}
    \item are all elements of this set also inside another set?
    \item does this set contain all elements of another set?
    \item do these two sets share no elements at all?
\end{itemize}

These methods return Boolean values. They answer with \texttt{True} or \texttt{False}.

\subsubsection{Containment and Scope Testing: Identifying Subsets \texttt{.issubset()} and Supersets \texttt{.issuperset()}}

A set \texttt{a} is a subset of set \texttt{b} when every element of \texttt{a} also appears in \texttt{b}.

\begin{verbatim}
small = {"spam", "eggs"}
large = {"spam", "eggs", "larch", 42}

print(f"small: {small!r}")
print(f"large: {large!r}")
print()

print(f"small.issubset(large): {small.issubset(large)}")
print(f"large.issubset(small): {large.issubset(small)}")
\end{verbatim}

Possible output:

\begin{verbatim}
small: {'eggs', 'spam'}
large: {42, 'eggs', 'larch', 'spam'}

small.issubset(large): True
large.issubset(small): False
\end{verbatim}

The first result is \texttt{True} because every element in \texttt{small} exists inside \texttt{large}. The second result is \texttt{False} because \texttt{large} has elements that \texttt{small} does not have.

A superset is the same relationship viewed from the opposite direction. A set \texttt{a} is a superset of set \texttt{b} when \texttt{a} contains every element of \texttt{b}.

\begin{verbatim}
small = {"spam", "eggs"}
large = {"spam", "eggs", "larch", 42}

print(f"large.issuperset(small): {large.issuperset(small)}")
print(f"small.issuperset(large): {small.issuperset(large)}")
\end{verbatim}

Possible output:

\begin{verbatim}
large.issuperset(small): True
small.issuperset(large): False
\end{verbatim}

The two method calls express the same relationship from opposite sides:

\begin{verbatim}
small.issubset(large)
large.issuperset(small)
\end{verbatim}

Both mean that \texttt{large} covers all elements from \texttt{small}.

The operators \texttt{<=} and \texttt{>=} can also be used.

\begin{verbatim}
small = {"Arthur", "Brian"}
large = {"Arthur", "Brian", "Lancelot"}

print(f"small <= large: {small <= large}")
print(f"large >= small: {large >= small}")
\end{verbatim}

Possible output:

\begin{verbatim}
small <= large: True
large >= small: True
\end{verbatim}

The method names are usually clearer in teaching code because they say the relationship directly.

Subset does not mean smaller by length only. It means that all members are included.

\begin{verbatim}
a = {"Homer", "Marge"}
b = {"Lisa", "Bart"}

print(f"len(a): {len(a)}")
print(f"len(b): {len(b)}")
print(f"a.issubset(b): {a.issubset(b)}")
\end{verbatim}

Possible output:

\begin{verbatim}
len(a): 2
len(b): 2
a.issubset(b): False
\end{verbatim}

The two sets have the same length, but neither contains the other's elements.

An empty set is a subset of every set, because there is no element inside the empty set that violates the containment rule.

\begin{verbatim}
empty = set()
names = {"Jerry", "Elaine", "George"}

print(f"empty: {empty!r}")
print(f"names: {names!r}")
print()

print(f"empty.issubset(names): {empty.issubset(names)}")
print(f"names.issuperset(empty): {names.issuperset(empty)}")
\end{verbatim}

Possible output:

\begin{verbatim}
empty: set()
names: {'Elaine', 'George', 'Jerry'}

empty.issubset(names): True
names.issuperset(empty): True
\end{verbatim}

This can feel strange at first, but the rule is literal: every element of \texttt{empty} is also in \texttt{names}. Since \texttt{empty} has no elements, there is no counterexample.

The methods can also accept other iterable objects.

\begin{verbatim}
allowed = {"Homer", "Marge", "Lisa", "Bart"}
requested = ["Lisa", "Bart"]

print(f"allowed: {allowed!r}")
print(f"requested: {requested!r}")
print()

print(f"allowed.issuperset(requested): {allowed.issuperset(requested)}")
\end{verbatim}

Possible output:

\begin{verbatim}
allowed: {'Lisa', 'Homer', 'Bart', 'Marge'}
requested: ['Lisa', 'Bart'}

allowed.issuperset(requested): True
\end{verbatim}

The list \texttt{requested} is traversed as a source of values. The relationship question is still set-like: does \texttt{allowed} contain all requested values?

A practical use is permission checking.

\begin{verbatim}
user_permissions = {"read", "write", "export"}
required_permissions = {"read", "export"}

if user_permissions.issuperset(required_permissions):
    print("operation allowed")
else:
    print("operation blocked")
\end{verbatim}

Possible output:

\begin{verbatim}
operation allowed
\end{verbatim}

The code asks whether the user's permission domain covers the required permission domain.

\subsubsection{Intersection Disjoint Verification via \texttt{.isdisjoint()}}

Two sets are disjoint when they share no elements. The method \texttt{.isdisjoint()} returns \texttt{True} if the intersection would be empty.

\begin{verbatim}
python_words = {"spam", "eggs", "larch"}
simpsons_words = {"D'oh!", "Mmm... donuts", "excellent"}

print(f"python_words:   {python_words!r}")
print(f"simpsons_words: {simpsons_words!r}")
print()

print(f"python_words.isdisjoint(simpsons_words): "
      f"{python_words.isdisjoint(simpsons_words)}")
\end{verbatim}

Possible output:

\begin{verbatim}
python_words:   {'eggs', 'larch', 'spam'}
simpsons_words: {'excellent', "D'oh!", 'Mmm... donuts'}

python_words.isdisjoint(simpsons_words): True
\end{verbatim}

The result is \texttt{True} because the two sets do not share any value.

If even one value is shared, the sets are not disjoint.

\begin{verbatim}
left = {"spam", "eggs", "larch"}
right = {"D'oh!", "spam", "Serenity now!"}

print(f"left:  {left!r}")
print(f"right: {right!r}")
print()

print(f"left.isdisjoint(right): {left.isdisjoint(right)}")
print(f"left & right: {left & right}")
\end{verbatim}

Possible output:

\begin{verbatim}
left:  {'eggs', 'larch', 'spam'}
right: {"D'oh!", 'Serenity now!', 'spam'}

left.isdisjoint(right): False
left & right: {'spam'}
\end{verbatim}

The intersection contains \texttt{"spam"}, so the sets are not disjoint.

The method \texttt{.isdisjoint()} is often clearer and more direct than computing the intersection manually.

\begin{verbatim}
blocked_users = {"Newman", "Burns", "Sideshow Bob"}
current_users = {"Jerry", "Elaine", "George"}

print(f"blocked_users: {blocked_users!r}")
print(f"current_users: {current_users!r}")
print()

if blocked_users.isdisjoint(current_users):
    print("no blocked users are active")
else:
    print("at least one blocked user is active")
\end{verbatim}

Possible output:

\begin{verbatim}
blocked_users: {'Sideshow Bob', 'Burns', 'Newman'}
current_users: {'Elaine', 'George', 'Jerry'}

no blocked users are active
\end{verbatim}

This method also accepts iterable objects.

\begin{verbatim}
blocked = {"Newman", "Burns", "Sideshow Bob"}
incoming = ["Jerry", "Elaine", "Newman"]

print(f"blocked: {blocked!r}")
print(f"incoming: {incoming!r}")
print()

print(f"blocked.isdisjoint(incoming): {blocked.isdisjoint(incoming)}")
\end{verbatim}

Possible output:

\begin{verbatim}
blocked: {'Sideshow Bob', 'Burns', 'Newman'}
incoming: ['Jerry', 'Elaine', 'Newman']

blocked.isdisjoint(incoming): False
\end{verbatim}

The list contains \texttt{"Newman"}, which is also in the blocked set. Therefore, the domains are not disjoint.

The relationship methods can be summarized as follows:

\begin{center}
\begin{tabular}{lll}
\textbf{Method} & \textbf{Question} & \textbf{Returns} \\
\hline
\texttt{a.issubset(b)} & Are all elements of \texttt{a} in \texttt{b}? & \texttt{bool} \\
\texttt{a.issuperset(b)} & Does \texttt{a} contain all elements of \texttt{b}? & \texttt{bool} \\
\texttt{a.isdisjoint(b)} & Do \texttt{a} and \texttt{b} share no elements? & \texttt{bool} \\
\end{tabular}
\end{center}

These methods do not change either set.

\begin{verbatim}
a = {"spam", "eggs"}
b = {"spam", "eggs", "larch"}

print(f"a before: {a!r}")
print(f"b before: {b!r}")

result = a.issubset(b)

print()
print(f"result: {result}")
print(f"a after:  {a!r}")
print(f"b after:  {b!r}")
\end{verbatim}

Possible output:

\begin{verbatim}
a before: {'eggs', 'spam'}
b before: {'eggs', 'spam', 'larch'}

result: True
a after:  {'eggs', 'spam'}
b after:  {'eggs', 'spam', 'larch'}
\end{verbatim}

They are inspection operations, not mutation operations. They read the membership domains and answer a structural question about them.